\documentclass[12pt]{article}
\usepackage{amsmath,amssymb,amsthm}

\newtheorem{theorem}{Theorem}

\begin{document}

\begin{theorem}
The Diophantine equation
\[
6 + x^3 + x y^2 z + z^2 = 0
\]
has no integer solutions.
\end{theorem}

\begin{proof}
Suppose that $(x,y,z)\in\mathbb{Z}^3$ satisfies
\[
6 + x^3 + x y^2 z + z^2 = 0.
\tag{$*$}
\]

\medskip

\noindent\textbf{Step 1: Parity of $x$.}
Reducing modulo $4$, if $x$ were even then $x^3$ and $x y^2 z$ would be even, hence
\[
z^2 \equiv -6 \equiv 2 \pmod{4},
\]
which is impossible since a square modulo $4$ is $0$ or $1$. Thus $x$ is odd.

\medskip

\noindent\textbf{Step 2: Parity of $y$ and $z$.}
Reducing modulo $2$, since $x$ is odd we have $x^3 \equiv 1 \pmod{2}$, so
\[
1 + y^2 z + z \equiv 0 \pmod{2},
\]
i.e.\ $z(y^2+1) \equiv 1 \pmod{2}$. Hence $z$ is odd and $y^2+1$ is odd, which forces $y$
to be even.

\medskip

\noindent\textbf{Step 3: Excluding $3 \mid x$.}
Suppose $3 \mid x$. Then reducing $(*)$ modulo $3$ gives
\[
z^2 \equiv 0 \pmod{3},
\]
so $3 \mid z$. Write $x=3a$ and $z=3b$. Substituting into $(*)$ yields
\[
6 + 27a^3 + 9a y^2 b + 9b^2 = 0.
\]
Reducing modulo $9$ gives $6 \equiv 0 \pmod{9}$, a contradiction. Hence $3 \nmid x$.

\medskip

\noindent\textbf{Step 4: Coprimality of $x$ and $z$.}
Let $d=\gcd(x,z)$. Then $d \mid x^3$, $d \mid z^2$, and $d \mid x y^2 z$, so from $(*)$
\[
d \mid 6.
\]
Since $x$ and $z$ are both odd, $d$ is odd, hence $d \in \{1,3\}$. By Step 3, $3 \nmid x$,
so $d=1$. Thus $\gcd(x,z)=1$.

\medskip

\noindent\textbf{Step 5: Bounding $x$.}
Rearranging $(*)$ gives
\[
x^3 + 6 = -z(z + xy^2).
\]
Since $x \mid x^3$ and $x \mid xy^2$, we have $x \mid z(z + xy^2)$.  Because
$\gcd(x,z)=1$, it follows that $x \mid (z+xy^2)$.  But $z+xy^2 \equiv z \pmod{x}$,
so $x \mid z$.  Since $\gcd(x,z)=1$, we conclude $|x| = 1$, i.e.\ $x = \pm 1$.

\medskip

\noindent\textbf{Step 6: Final contradiction.}

\smallskip

\emph{Case $x=1$.} The equation becomes
\[
z^2 + y^2 z + 7 = 0.
\]
From Step 5's factored form, $z \mid (1+6)=7$, so $z \in \{\pm 1,\pm 7\}$.
Solving $y^2 = -(z^2+7)/z$ for each value:
\[
z = -1:\; y^2 = 8,\qquad z = -7:\; y^2 = 8,\qquad z = 1:\; y^2 = -8,\qquad z = 7:\; y^2 = -8.
\]
The first two are impossible since $8$ is not a perfect square; the last two are impossible
since $y^2 \geq 0$.

\smallskip

\emph{Case $x=-1$.} The equation becomes
\[
z^2 - y^2 z + 5 = 0.
\]
From Step 5's factored form, $z \mid (-1+6)=5$, so $z \in \{\pm 1,\pm 5\}$.
Solving $y^2 = (z^2+5)/z$ for each value:
\[
z = 1:\; y^2 = 6,\qquad z = 5:\; y^2 = 6,\qquad z = -1:\; y^2 = -6,\qquad z = -5:\; y^2 = -6.
\]
The first two are impossible since $6$ is not a perfect square; the last two are impossible
since $y^2 \geq 0$.

\medskip

In all cases we obtain a contradiction. Therefore there are no integer solutions.
\end{proof}

\end{document}
